\section{Introduction}
\label{sec:introduction}

Behavior cloning (BC) turns robot learning into supervised learning from
expert state--action pairs. Its simplicity is attractive for manipulation,
where online exploration can be unsafe or expensive. Yet a policy trained on
expert states is evaluated on states induced by its own actions. Small errors
can therefore compound and move the robot outside the demonstration
distribution, a classical limitation of offline imitation
\citep{pomerleau1989alvinn,ross2011dagger}. How demonstrations are sampled
during optimization may influence both learning speed and the final policy.

Curriculum learning proposes to organize examples rather than sample them
uniformly \citep{bengio2009curriculum}. In manipulation, however, ``easy'' is
ambiguous. A task can be geometrically demanding; an expert can require a long
or circuitous execution; or a learner can predict a demonstration poorly.
These notions need not agree. Moreover, easy-to-hard is not the only plausible
direction: difficult examples can expose useful decision boundaries early,
motivating a reverse ordering related to hard-example and reverse-curriculum
ideas \citep{florensa2017reverse}.

We study three questions: \textbf{RQ1 (ranking agreement):} do Task-, Expert-,
and Learner-based scores produce consistent trajectory rankings?
\textbf{RQ2 (failure prediction):} do these scores predict closed-loop BC
failures? \textbf{RQ3 (curriculum effectiveness):} under matched training
budgets and initialization protocols, how do staged curricula based on these
scores compare with uniform sampling? The stages change cumulative exposure as
well as order, so RQ3 concerns demonstration-pool sampling schedules rather
than pure permutations with equal exposure.

We construct all scores on one frozen split, evaluate five schedules, analyze
failure prediction and ranking corruption, and repeat the comparison on
stacking as an exploratory secondary evaluation. We propose no new curriculum
algorithm; the contribution is a controlled empirical analysis.

The results are mixed in an informative way. On pick-and-place, task and
expert difficulty correlate at $0.842$, while learner difficulty correlates
negatively with each. Learner difficulty has the largest point-estimate
failure AUROC, but its confidence interval includes chance. Reverse Learner
has the highest observed pick-and-place Fresh-ID mean, although its seed-level
advantage is inconsistent. Stacking is near the performance floor, where no
curriculum outperforms Uniform. Thus the evidence cautions against treating
either difficulty definitions or schedule effects as task-independent.
